\section{Module \texttt{display.c}}

Le module \texttt{display.c} prend en charge les widgets d'affichage pur : barre de progression, spinner, label et image. Son rôle est d'exposer des constructeurs simples au-dessus des widgets GTK standards.

\subsection{Fonction \texttt{create\_progress\_bar}}

\begin{lstlisting}[language=C]
GtkWidget *create_progress_bar(const progress_bar_config *config) {
    GtkWidget *progress = gtk_progress_bar_new();

    if (config == NULL) {
        return progress;
    }

    gtk_progress_bar_set_fraction(GTK_PROGRESS_BAR(progress), config->fraction);
    gtk_progress_bar_set_show_text(GTK_PROGRESS_BAR(progress), config->show_text);
    gtk_progress_bar_set_inverted(GTK_PROGRESS_BAR(progress), config->inverted);
    gtk_orientable_set_orientation(GTK_ORIENTABLE(progress), config->orientation);
    if (config->text != NULL) {
        gtk_progress_bar_set_text(GTK_PROGRESS_BAR(progress), config->text);
    }
    if (config->pulse_once) {
        gtk_progress_bar_pulse(GTK_PROGRESS_BAR(progress));
    }
    apply_widget_style(progress, &config->style);

    return progress;
}
\end{lstlisting}

\subsection{Fonction \texttt{create\_spinner}}

\begin{lstlisting}[language=C]
GtkWidget *create_spinner(const spinner_config *config) {
    GtkWidget *spinner = gtk_spinner_new();

    if (config != NULL) {
        if (config->size > 0) {
            gtk_widget_set_size_request(spinner, config->size, config->size);
        }
        if (config->spinning) {
            gtk_spinner_start(GTK_SPINNER(spinner));
        }
        apply_widget_style(spinner, &config->style);
    }

    return spinner;
}
\end{lstlisting}

\subsection{Fonction \texttt{create\_label}}

\begin{lstlisting}[language=C]
GtkWidget *create_label(const label_config *config) {
    GtkWidget *label = gtk_label_new(config && config->text ? config->text : "");

    if (config == NULL) {
        return label;
    }

    gtk_label_set_selectable(GTK_LABEL(label), config->selectable);
    gtk_label_set_wrap(GTK_LABEL(label), config->wrap);
    gtk_label_set_ellipsize(GTK_LABEL(label), config->ellipsize);
    if (config->max_width_chars > 0) {
        gtk_label_set_max_width_chars(GTK_LABEL(label), config->max_width_chars);
    }
    apply_widget_style(label, &config->style);

    return label;
}
\end{lstlisting}

\subsection{Fonction \texttt{create\_image}}

\begin{lstlisting}[language=C]
GtkWidget *create_image(const image_config *config) {
    GtkWidget *widget;

    if (config == NULL) {
        return gtk_image_new();
    }

    if (config->file_path != NULL && config->file_path[0] != '\0') {
        GtkWidget *pic = gtk_picture_new_for_filename(config->file_path);
        gtk_picture_set_can_shrink(GTK_PICTURE(pic), TRUE);
        gtk_picture_set_content_fit(
            GTK_PICTURE(pic),
            config->keep_aspect ? GTK_CONTENT_FIT_CONTAIN : GTK_CONTENT_FIT_FILL
        );
        widget = pic;
    } else if (config->icon_name != NULL && config->icon_name[0] != '\0') {
        GtkWidget *img = gtk_image_new_from_icon_name(config->icon_name);
        if (config->pixel_size > 0) {
            gtk_image_set_pixel_size(GTK_IMAGE(img), config->pixel_size);
        }
        widget = img;
    } else {
        widget = gtk_image_new();
    }

    apply_widget_style(widget, &config->style);
    return widget;
}
\end{lstlisting}

\subsection{APIs GTK utilisées}

\begin{itemize}
    \item \textbf{\texttt{gtk\_progress\_bar\_set\_fraction()}} : Définit la fraction de remplissage de la barre de progression, acceptant une valeur flottante entre 0.0 et 1.0. Elle est utilisée dans \texttt{create\_progress\_bar} pour initialiser le niveau d'avancement visuel.
    \item \textbf{\texttt{gtk\_progress\_bar\_pulse()}} : Fait basculer la barre de progression en mode "activité indéterminée", déplaçant un bloc de manière cyclique. Le wrapper l'appelle pour indiquer que le programme est en cours de traitement sans pouvoir estimer le temps restant.
    \item \textbf{\texttt{gtk\_spinner\_start()}} : Démarre l'animation de rotation du widget spinner. Le module l'utilise pour activer visuellement l'indicateur de chargement dès sa création si la configuration le spécifie.
    \item \textbf{\texttt{gtk\_label\_set\_ellipsize()}} : Détermine le mode de troncature du texte (début, milieu ou fin) par des points de suspension lorsque l'espace est insuffisant. Elle permet au wrapper de garantir une mise en page propre même avec des labels imprévus.
    \item \textbf{\texttt{gtk\_picture\_new\_for\_filename()}} : Crée un widget d'image haute performance à partir d'un fichier local. Contrairement à \texttt{GtkImage}, cette API utilisée par le module facilite le redimensionnement proportionnel des images chargées.
    \item \textbf{\texttt{gtk\_image\_new\_from\_icon\_name()}} : Instancie un widget d'image affichant une icône issue du thème de l'utilisateur. Le wrapper l'exploite pour intégrer des symboles graphiques standards dans l'interface de manière simple et légère.
\end{itemize}

\newpage
